\documentclass[10pt]{article}
\usepackage[T1]{fontenc}
\usepackage[utf8]{inputenc}
\usepackage[french]{babel}
\usepackage{amsmath,amssymb,mathrsfs,mathpazo}
\usepackage{geometry}
\geometry{a4paper,margin=2cm}
\usepackage{graphicx,array,tabularx,tabularray,colortbl,multirow,multicol,booktabs,enumitem,pifont,ulem,fancybox,diagbox,cancel}
\usepackage{pstricks,pstricks-add,pst-plot,pst-tree,pst-node,pst-text,pst-func,pst-math,pst-bspline,pst-3dplot,pst-eucl,pst-solides3d}
\let\pscalculate\undefined
\usepackage{pst-bezier}
\usepackage{tikz,pgfplots,tkz-tab}
\usetikzlibrary{arrows,arrows.meta,patterns,calc,positioning,trees,decorations.markings,shapes.geometric,shapes.misc,shapes.arrows,intersections,decorations.pathmorphing,decorations.pathreplacing,backgrounds,shadows,babel}
\pgfplotsset{compat=1.18}
\usepgfplotslibrary{fillbetween}
\usepackage[np]{numprint}
\usepackage{listings,makecell,pas-tableur,twoopt}
\let\barre\undefined
\usepackage[pstricks]{tablvar}
\definecolor{darkgreen}{RGB}{0,100,0}
\providecommand{\red}{\color{red}}
\providecommand{\blue}{\color{blue}}
\providecommand{\green}{\color{green}}
\providecommand{\black}{\color{black}}
\providecommand{\white}{\color{white}}
\providecommand{\gray}{\color{gray}}
\providecommand{\cyan}{\color{cyan}}
\providecommand{\magenta}{\color{magenta}}
\providecommand{\yellow}{\color{yellow}}
\providecommand{\violet}{\color{violet}}
\providecommand{\orange}{\color{orange}}
\providecommand{\euro}{\texteuro}
\providecommand{\EUR}{\texteuro}
\providecommand{\ieme}{\textsuperscript{e}}
\providecommand{\degres}{\ensuremath{^\circ}}
\providecommand{\degree}{\ensuremath{^\circ}}
\providecommand{\vv}[1]{\overrightarrow{#1}}
\providecommand{\Coord}[1]{#1}
\providecommand{\decosix}{}
\providecommand{\starredbullet}{\ensuremath{\bullet}}
\pagestyle{empty}
\setlength{\parindent}{0pt}
\usepackage{adjustbox}\DeclareUnicodeCharacter{2236}{:}\DeclareUnicodeCharacter{2190}{\ensuremath{\leftarrow}}
\providecommand{\R}{\mathbb{R}}
\providecommand{\N}{\mathbb{N}}
\providecommand{\Z}{\mathbb{Z}}
\providecommand{\Q}{\mathbb{Q}}
\providecommand{\D}{\mathbb{D}}
\providecommand{\C}{\mathbb{C}}
\providecommand{\e}{{\mathrm{e}}}
\providecommand{\vect}[1]{\overrightarrow{#1}}
\providecommand{\vectt}[1]{\overrightarrow{\mathrm{#1}}}
\providecommand{\vvt}[1]{\overrightarrow{\mathrm{#1}}}
\providecommand{\np}[2][]{{#2~#1}}
\providecommand{\vv}[1]{{\overrightarrow{#1}}}
\renewcommand{\d}{{\mathrm{d}}}
\providecommand{\degres}{{^\circ}}
\providecommand{\degre}{{^\circ}}
\providecommand{\blue}{\color{blue}}
\providecommand{\red}{\color{red}}
\providecommand{\enskip}{\quad}
\providecommand{\llbracket}{\lbrack\!\lbrack}
\providecommand{\rrbracket}{\rbrack\!\rbrack}
\providecommand{\euro}{{\text{€}}}
\providecommand{\no}{{\text{n°}}}
\providecommand{\interval}[2]{{[#1;#2]}}
\providecommand{\intervFO}[2]{{[#1;#2[}}
\providecommand{\dotfill}{\ldots}

\providecommand{\ds}{\displaystyle}
\providecommand{\textding}[1]{#1}
\providecommand{\Oij}{\(\mathrm{O};\vec{\imath},\vec{\jmath}\)}
\providecommand{\Oijk}{\(\mathrm{O};\vec{\imath},\vec{\jmath},\vec{k}\)}
\providecommand{\Ouv}{\(\mathrm{O};\vec{u},\vec{v}\)}
\providecommand{\intervFF}[2]{\ensuremath{[#1;#2]}}
\renewcommand{\euro}{\ifmmode\text{\texteuro}\else\texteuro\fi}
\DeclareUnicodeCharacter{00B0}{\ifmmode^{\circ}\else\textdegree\fi}
\DeclareUnicodeCharacter{20AC}{\ifmmode\text{\texteuro}\else\texteuro\fi}
% Énoncés : Éducation nationale. Transcription et corrigés : APMEP, Denis Vergès et contributeurs. Découpage et classement : CoopMaths, Sébastien Lozano, Éric Elter et contributeurs. Adaptation de la présentation : Maths Explorer.
% CC BY-SA : https://coopmaths.fr/www/about/
% https://coopmaths.fr/alea/static/eam/2026/tex/eam_specifique_2026_06_ameriquenord_3_cor.tex
\begin{document}
\noindent\textit{APMEP / CoopMaths. CC BY-SA. Présentation adaptée par Maths Explorer.}\par\medskip
On considère ci-dessous la représentation graphique d'une fonction $f$ définie sur $[-2~;~4]$. On a également tracé sa tangente au point $A$ d'abscisse $-1$.
	
	\medskip
	
	\begin{minipage}{0.46\linewidth}
		\begin{enumerate}
			\item Par lecture graphique, donner la valeur de :
			\begin{enumerate}
				\item $f(3)$
				
				\begin{color}{green!50!black}
				$f(3)\approx 5$
		 		\end{color}

				
				
				\item $f'(-1)$
				
				\begin{color}{green!50!black}
				$f'(-1)$ est donné par le coefficient directeur de la tangente à la courbe représentant $f$ au point d'abscisse $-1$.				
				
				$f'(-1)\approx 4$
		 		\end{color}
				
			\end{enumerate}
		\end{enumerate}
		
		
		
	\end{minipage}
	\hfill
	\begin{minipage}{0.5\linewidth}
		\centering
		\begin{tikzpicture}
			\begin{axis}[/pgf/number format/.cd, use comma,
				scaled ticks=false,
				x=8mm, y=8mm,
				xmin=-2.6, xmax=4.6, ymin=-0.6, ymax=9.8,
				axis lines=center,
				xtick={-3,-2,-1,0,1,2,3,4,5}, ytick={-1,0,1,2,3,4,5,6,7,8,9,10},
				tick label style={font=\scriptsize},
				minor tick num=1,
				grid=both,
				grid style={gray!70,dotted},
				clip=true]
				\addplot[teal,line width=1.5pt,samples=100,domain=-2:4]{-x^2+2*x+8};
				\fill[teal] (axis cs:-2,0) circle (1.5pt);
				\fill[teal] (axis cs: 4,0) circle (1.5pt);
				\addplot[red,line width=1.5pt,domain=-2.4:0.3]{4*x+9};
				\node[below left,font=\scriptsize] at (axis cs:0,0) {0};
				\node[above left,font=\footnotesize] at (axis cs:-1,5) {$A$};
				\node[font=\footnotesize,above left] at (axis cs:4.6,0) {$x$};
				\node[font=\footnotesize,below left] at (axis cs:0,9.8) {$y$};
				\filldraw (axis cs:-1,5) circle (1.5pt);
				\draw[orange,preaction={fill,opacity=0.5,fill=orange}] (axis cs:-1,5) -- (axis cs:0,9) -- (axis cs:0,5) node[midway, right] {\textbf{4}} -- cycle node[midway, below] {1} ;
				\draw[dashed, purple, opacity=0.5,line width=1.5pt] (axis cs:0,5) -- (axis cs:3,5) --  (axis cs:3,0) ;
				\node[rotate=45, purple,line width=1.5pt,draw,cross out,inner sep=0pt,minimum width=5pt,minimum height=5pt, opacity=0.7] at (axis cs:3,5) {};
				\node[purple,above right] at (axis cs:3,5) {$(3,\mathbf{5})$};
			\end{axis}
		\end{tikzpicture}
	\end{minipage}
	
	\begin{enumerate}[start=2]
			\item On admet que la fonction $f$ est définie sur $[-2~;~4]$ par \quad $f(x)=-x^{2}+2x+8$.
			\begin{enumerate}
				\item Calculer $f'(x)$ pour $x$ appartenant à $[-2~;~4]$.
				
				\begin{color}{green!50!black}
				$f'(x) = - (2\times x) + 2 \times 1 = -2x+2$
		 		\end{color}
				
			
				\item Étudier le signe de $f'(x)$ sur $[-2~;~4]$.
				
				\begin{color}{green!50!black}
				\[
				\begin{array}{ccc}
				-2x+2 = 0	 & \iff & 2x = 2 \\
						 & \iff & x = 1 \\
				\end{array}
				\]
				
				La fonction $f'$ est la restriction d'une fonction affine sur $[-2~;~4]$. Le taux d'accroissement est négatif donc elle est décroissante. Elle admet pour racine le nombre $1$. Ainsi $f'(x)$ est positif si $x$ est inférieur à $1$ et négatif si $x$ est supérieur à $1$.
				
				
		 		\end{color}
				
				
			\end{enumerate}
			\item Donner les variations de $f$ sur $[-2~;~4]$.
			
			\begin{color}{green!50!black}
			$f(-2) = -(-2)^2 +2\times(-2) + 8 = -4 - 4 + 8 =0$
						
			$f(4) = -4^2 +2\times 4 + 8 = -16 + 8 + 8 =0$
			
			$f(1) = -1^2 +2\times 1 + 8 = -1 + 2 + 8 = 9$
			
			\begin{tikzpicture}
				\tkzTabInit[lgt=3, espcl=4]{$x$ / 1, {signe de $f'(x)$} / 1.5, {variations de $f(x)$} / 3.5}{$\np{-2}$, $1$ , $4$}
				\tkzTabLine{,+,z,-,}
				\tkzTabVar{-/$0$ , +/$9$ , -/$0$ }
			\end{tikzpicture}
			\end{color}
				
			
		\end{enumerate}
	
	
	\medskip

\end{document}
